\documentclass[10pt,a4paper]{amsart}
\usepackage{amsmath,amssymb,amsthm,graphicx}
% Retained minimal preamble for the Stacks Project chapter sites-cohomology.tex (Cohomology on Sites).
% Verbatim excerpt of the chapter preamble: Xy-pic setup (native preamble lines 14-18) and the native
% theorem declarations with the mathematical macros (native preamble lines 154-217). Copied mechanically;
% nothing added, removed or redefined. The wrapper class and the amsmath/amssymb/amsthm/graphicx base
% packages are supplied by the assembler, so the class/externaldocument part of the native preamble is omitted.
% For commutative diagrams we use Xy-pic
\usepackage[all]{xy}

% We use 2cell for 2-commutative diagrams.
\xyoption{2cell}
% Theorem environments.
%
\theoremstyle{plain}
\newtheorem{theorem}[subsection]{Theorem}
\newtheorem{proposition}[subsection]{Proposition}
\newtheorem{lemma}[subsection]{Lemma}

\theoremstyle{definition}
\newtheorem{definition}[subsection]{Definition}
\newtheorem{example}[subsection]{Example}
\newtheorem{exercise}[subsection]{Exercise}
\newtheorem{situation}[subsection]{Situation}

\theoremstyle{remark}
\newtheorem{remark}[subsection]{Remark}
\newtheorem{remarks}[subsection]{Remarks}

\numberwithin{equation}{subsection}

% Macros
%
\def\lim{\mathop{\mathrm{lim}}\nolimits}
\def\colim{\mathop{\mathrm{colim}}\nolimits}
\def\Spec{\mathop{\mathrm{Spec}}}
\def\Hom{\mathop{\mathrm{Hom}}\nolimits}
\def\Ext{\mathop{\mathrm{Ext}}\nolimits}
\def\SheafHom{\mathop{\mathcal{H}\!\mathit{om}}\nolimits}
\def\SheafExt{\mathop{\mathcal{E}\!\mathit{xt}}\nolimits}
\def\Sch{\mathit{Sch}}
\def\Mor{\mathop{\mathrm{Mor}}\nolimits}
\def\Ob{\mathop{\mathrm{Ob}}\nolimits}
\def\Sh{\mathop{\mathit{Sh}}\nolimits}
\def\NL{\mathop{N\!L}\nolimits}
\def\CH{\mathop{\mathrm{CH}}\nolimits}
\def\proetale{{pro\text{-}\acute{e}tale}}
\def\etale{{\acute{e}tale}}
\def\QCoh{\mathit{QCoh}}
\def\Ker{\mathop{\mathrm{Ker}}}
\def\Im{\mathop{\mathrm{Im}}}
\def\Coker{\mathop{\mathrm{Coker}}}
\def\Coim{\mathop{\mathrm{Coim}}}

% Boxtimes
%
\DeclareMathSymbol{\boxtimes}{\mathbin}{AMSa}{"02}

%
% Macros for moduli stacks/spaces
%
\def\QCohstack{\mathcal{QC}\!\mathit{oh}}
\def\Cohstack{\mathcal{C}\!\mathit{oh}}
\def\Spacesstack{\mathcal{S}\!\mathit{paces}}
\def\Quotfunctor{\mathrm{Quot}}
\def\Hilbfunctor{\mathrm{Hilb}}
\def\Curvesstack{\mathcal{C}\!\mathit{urves}}
\def\Polarizedstack{\mathcal{P}\!\mathit{olarized}}
\def\Complexesstack{\mathcal{C}\!\mathit{omplexes}}
% \Pic is the operator that assigns to X its picard group, usage \Pic(X)
% \Picardstack_{X/B} denotes the Picard stack of X over B
% \Picardfunctor_{X/B} denotes the Picard functor of X over B
\def\Pic{\mathop{\mathrm{Pic}}\nolimits}
\def\Picardstack{\mathcal{P}\!\mathit{ic}}
\def\Picardfunctor{\mathrm{Pic}}
\def\Deformationcategory{\mathcal{D}\!\mathit{ef}}

% --- appended 2026-10-12 (Chief): the project own definitions of the reference/slogan/history
% display environments, copied mechanically from the shared Stacks preamble lines 8-12
% (\usepackage{verbatim} plus the three \newenvironment lines, which the chapter preamble keeps
% verbatim from the project preamble). Nothing added, removed or redefined: these environments
% carry no mathematical content and the project itself typesets them as comments. Without these
% three lines a counted statement whose own text uses them fails with Environment slogan undefined.
\usepackage{verbatim}
\newenvironment{reference}{\comment}{\endcomment}
\newenvironment{slogan}{\comment}{\endcomment}
\newenvironment{history}{\comment}{\endcomment}
\usepackage{cleveref}
\crefname{originaldefinition}{original definition}{original definitions}
\Crefname{originaldefinition}{Original definition}{Original definitions}
\makeatletter
\newlabel{derived-definition-derived-limit}{{08TC}{source}{}{sourcepointer-1}{}}
\newlabel{derived-definition-derived-limit@cref}{{[originaldefinition][0][]08TC}{[1][0][]source}{}{}{}}
\makeatother
\crefname{originallemma}{original lemma}{original lemmas}
\Crefname{originallemma}{Original lemma}{Original lemmas}
\makeatletter
\newlabel{derived-lemma-products}{{07KC}{source}{}{sourcepointer-2}{}}
\newlabel{derived-lemma-products@cref}{{[originallemma][0][]07KC}{[1][0][]source}{}{}{}}
\makeatother
\crefname{originallemma}{original lemma}{original lemmas}
\Crefname{originallemma}{Original lemma}{Original lemmas}
\makeatletter
\newlabel{homology-lemma-product-abelian-groups-exact}{{060K}{source}{}{sourcepointer-3}{}}
\newlabel{homology-lemma-product-abelian-groups-exact@cref}{{[originallemma][0][]060K}{[1][0][]source}{}{}{}}
\makeatother
\crefname{originalreference}{original reference}{original references}
\Crefname{originalreference}{Original reference}{Original references}
\makeatletter
\newlabel{lemma-is-limit-dimension}{{0D64}{source}{}{sourcepointer-4}{}}
\newlabel{lemma-is-limit-dimension@cref}{{[originalreference][0][]0D64}{[1][0][]source}{}{}{}}
\makeatother
\crefname{originalreference}{original reference}{original references}
\Crefname{originalreference}{Original reference}{Original references}
\makeatletter
\newlabel{lemma-w-contractible}{{090K}{source}{}{sourcepointer-5}{}}
\newlabel{lemma-w-contractible@cref}{{[originalreference][0][]090K}{[1][0][]source}{}{}{}}
\makeatother
\crefname{originalsection}{original section}{original sections}
\Crefname{originalsection}{Original section}{Original sections}
\makeatletter
\newlabel{section-derived-limits}{{0940}{source}{}{sourcepointer-6}{}}
\newlabel{section-derived-limits@cref}{{[originalsection][0][]0940}{[1][0][]source}{}{}{}}
\makeatother
\crefname{originaldefinition}{original definition}{original definitions}
\Crefname{originaldefinition}{Original definition}{Original definitions}
\makeatletter
\newlabel{sites-definition-sheaves-injective-surjective}{{00WM}{source}{}{sourcepointer-7}{}}
\newlabel{sites-definition-sheaves-injective-surjective@cref}{{[originaldefinition][0][]00WM}{[1][0][]source}{}{}{}}
\makeatother
\title{A basis of weakly contractible objects controls exactness, products and derived limits of modules on a site}
\author{The Stacks Project Authors}
\date{stacks-project:sites-cohomology.tex:local-snapshot-20261011; GNU FDL 1.2 or later}
\begin{document}
\maketitle
\noindent\textit{Editorial note.} All delivered mathematical text is the author's own, copied from the pinned Stacks Project chapter file sites-cohomology.tex; no statement or proof step was rewritten, reordered or removed.\par\medskip
\noindent\textit{Editorial note.} The counted proof is the author's complete argument: the proof is written out in full for all seven parts: the criterion on the basis with local lifting, the limit-dimension input, the surjectivity computation, the product exactness via the abelian-group lemma, and the derivation of the two derived-limit descriptions from the defining triangle.\par\medskip
\noindent\textit{Editorial note.} The label proposition-enough-weakly-contractibles names the counted statement, and the proof environment immediately following it in the source is delivered verbatim as its proof.\par\medskip
\noindent\textit{Editorial note.} The counted statement is carried by the author's own proposition environment in the chapter source.\par\medskip
\noindent\textit{Editorial note.} Pre-open gate: python operations/agent-authoring/tools/precard\_gate.py --source sites-cohomology.tex --label proposition-enough-weakly-contractibles --start 13800 --end 13918 --local-id 3547 returned RESULT=OK (L3547 free, occupancy FREE).\par\medskip
\noindent\textit{Editorial note.} The statement and proof are copied from the pinned chapter file; the cited inputs are kept as external pointers with their canonical Stacks tags.\par\medskip
\noindent\textit{Allowed standard prerequisites (editorial).}\begin{itemize}
\item Weakly contractible objects. Same chapter, the lemma recording that sections over a weakly contractible object form an exact functor with vanishing higher cohomology. Supplies the exactness used in parts (1) and (4) and the surjectivity in part (3).
\item Limit dimension of the site. Same chapter, the lemma relating the vanishing of derived limits to the limit dimension d. Supplies part (2) with d = 0.
\item Exactness of finite products of abelian groups and products in the derived category. Homology, the lemma that a finite product of exact sequences of abelian groups is exact, and Derived Categories, the lemma describing products in the derived category, together with the definition of the derived limit. Supplies parts (4), (5) and the distinguished triangle used for parts (6) and (7).
\end{itemize}\medskip
\ifdefined\hypertarget\hypertarget{sourcepointer-1}{}\fi
\noindent\textit{Source reference.} Original definition 08TC: sources/stacks-chow/tags.txt, source lines 11160--11160. Prerequisite pointer invoked at proof line(s) 13899; it names a standard fact of a sibling chapter that the counted argument uses, and every step of the counted proof built on it is retained verbatim in the delivered body.\par\medskip
\ifdefined\hypertarget\hypertarget{sourcepointer-2}{}\fi
\noindent\textit{Source reference.} Original lemma 07KC: sources/stacks-chow/tags.txt, source lines 9652--9652. Prerequisite pointer invoked at proof line(s) 13894; it names a standard fact of a sibling chapter that the counted argument uses, and every step of the counted proof built on it is retained verbatim in the delivered body.\par\medskip
\ifdefined\hypertarget\hypertarget{sourcepointer-3}{}\fi
\noindent\textit{Source reference.} Original lemma 060K: sources/stacks-chow/tags.txt, source lines 7729--7729. Prerequisite pointer invoked at proof line(s) 13889; it names a standard fact of a sibling chapter that the counted argument uses, and every step of the counted proof built on it is retained verbatim in the delivered body.\par\medskip
\ifdefined\hypertarget\hypertarget{sourcepointer-4}{}\fi
\noindent\textit{Source reference.} Original reference 0D64: sources/stacks-sites-cohomology/pinned-original-source.tex, source lines 5440--5440. Prerequisite pointer invoked at proof line(s) 13854; it names a standard fact of the same chapter that the counted argument uses, and every step of the counted proof built on it is retained verbatim in the delivered body.\par\medskip
\ifdefined\hypertarget\hypertarget{sourcepointer-5}{}\fi
\noindent\textit{Source reference.} Original reference 090K: sources/stacks-sites-cohomology/pinned-original-source.tex, source lines 13773--13773. Prerequisite pointer invoked at proof line(s) 13839; it names a standard fact of the same chapter that the counted argument uses, and every step of the counted proof built on it is retained verbatim in the delivered body.\par\medskip
\ifdefined\hypertarget\hypertarget{sourcepointer-6}{}\fi
\noindent\textit{Source reference.} Original section 0940: sources/stacks-sites-cohomology/pinned-original-source.tex, source lines 5016--5016. Prerequisite pointer invoked at proof line(s) 13897; it names a standard fact of the same chapter that the counted argument uses, and every step of the counted proof built on it is retained verbatim in the delivered body.\par\medskip
\ifdefined\hypertarget\hypertarget{sourcepointer-7}{}\fi
\noindent\textit{Source reference.} Original definition 00WM: sources/stacks-chow/tags.txt, source lines 1188--1188. Prerequisite pointer invoked at proof line(s) 13863; it names a standard fact of a sibling chapter that the counted argument uses, and every step of the counted proof built on it is retained verbatim in the delivered body.\par\medskip
\begin{proposition}
\label{proposition-enough-weakly-contractibles}
Let $\mathcal{C}$ be a site. Let $\mathcal{B} \subset \Ob(\mathcal{C})$
such that every $U \in \mathcal{B}$ is weakly contractible and
every object of $\mathcal{C}$ has a covering by elements of $\mathcal{B}$.
Let $\mathcal{O}$ be a sheaf of rings on $\mathcal{C}$. Then
\begin{enumerate}
\item A complex $\mathcal{F}_1 \to \mathcal{F}_2 \to \mathcal{F}_3$
of $\mathcal{O}$-modules is exact, if and only if
$\mathcal{F}_1(U) \to \mathcal{F}_2(U) \to \mathcal{F}_3(U)$
is exact for all $U \in \mathcal{B}$.
\item Every object $K$ of $D(\mathcal{O})$ is a derived limit
of its canonical truncations: $K = R\lim \tau_{\geq -n} K$.
\item Given an inverse system
$\ldots \to \mathcal{F}_3 \to \mathcal{F}_2 \to \mathcal{F}_1$
with surjective transition maps, the projection
$\lim \mathcal{F}_n \to \mathcal{F}_1$ is surjective.
\item Products are exact on $\textit{Mod}(\mathcal{O})$.
\item Products on $D(\mathcal{O})$ can be computed by taking
products of any representative complexes.
\item If $(\mathcal{F}_n)$ is an inverse system of $\mathcal{O}$-modules,
then $R^p\lim \mathcal{F}_n = 0$ for all $p > 1$ and
$$
R^1\lim \mathcal{F}_n  =
\Coker(\prod \mathcal{F}_n \to \prod \mathcal{F}_n)
$$
where the map is $(x_n) \mapsto (x_n - f(x_{n + 1}))$.
\item If $(K_n)$ is an inverse system of objects of $D(\mathcal{O})$,
then there are short exact sequences
$$
0 \to R^1\lim H^{p - 1}(K_n) \to H^p(R\lim K_n) \to
\lim H^p(K_n) \to 0
$$
\end{enumerate}
\end{proposition}

\begin{proof}
Proof of (1). If the sequence is exact, then evaluating at any
weakly contractible element of $\mathcal{C}$ gives an exact
sequence by Lemma \ref{lemma-w-contractible}. Conversely, assume that
$\mathcal{F}_1(U) \to \mathcal{F}_2(U) \to \mathcal{F}_3(U)$
is exact for all $U \in \mathcal{B}$.
Let $V$ be an object of $\mathcal{C}$ and let
$s \in \mathcal{F}_2(V)$ be an element of the kernel of
$\mathcal{F}_2 \to \mathcal{F}_3$. By assumption there exists
a covering $\{U_i \to V\}$ with $U_i \in \mathcal{B}$.
Then $s|_{U_i}$ lifts to a section $s_i \in \mathcal{F}_1(U_i)$.
Thus $s$ is a section of the image sheaf
$\Im(\mathcal{F}_1 \to \mathcal{F}_2)$.
In other words, the sequence
$\mathcal{F}_1 \to \mathcal{F}_2 \to \mathcal{F}_3$
is exact.

\medskip\noindent
Proof of (2). This holds by Lemma \ref{lemma-is-limit-dimension} with $d = 0$.

\medskip\noindent
Proof of (3). Let $(\mathcal{F}_n)$ be a system as in (2) and set
$\mathcal{F} = \lim \mathcal{F}_n$. If $U \in \mathcal{B}$, then
$\mathcal{F}(U) = \lim \mathcal{F}_n(U)$
surjects onto $\mathcal{F}_1(U)$ as all the transition maps
$\mathcal{F}_{n + 1}(U) \to \mathcal{F}_n(U)$ are surjective.
Thus $\mathcal{F} \to \mathcal{F}_1$ is surjective by
Sites, Definition \ref{sites-definition-sheaves-injective-surjective}
and the assumption that every object
has a covering by elements of $\mathcal{B}$.

\medskip\noindent
Proof of (4). Let
$\mathcal{F}_{i, 1} \to \mathcal{F}_{i, 2} \to \mathcal{F}_{i, 3}$
be a family of exact sequences of $\mathcal{O}$-modules.
We want to show that
$\prod \mathcal{F}_{i, 1} \to \prod \mathcal{F}_{i, 2} \to
\prod \mathcal{F}_{i, 3}$ is exact. We use the criterion of (1).
Let $U \in \mathcal{B}$. Then
$$
(\prod \mathcal{F}_{i, 1})(U) \to
(\prod \mathcal{F}_{i, 2})(U) \to
(\prod \mathcal{F}_{i, 3})(U)
$$
is the same as
$$
\prod \mathcal{F}_{i, 1}(U) \to
\prod \mathcal{F}_{i, 2}(U) \to
\prod \mathcal{F}_{i, 3}(U)
$$
Each of the sequences
$\mathcal{F}_{i, 1}(U) \to \mathcal{F}_{i, 2}(U) \to \mathcal{F}_{i, 3}(U)$
are exact by (1). Thus the displayed sequences are exact by
Homology, Lemma \ref{homology-lemma-product-abelian-groups-exact}.
We conclude by (1) again.

\medskip\noindent
Proof of (5). Follows from (4) and (slightly generalized)
Derived Categories, Lemma \ref{derived-lemma-products}.

\medskip\noindent
Proof of (6) and (7). We refer to Section \ref{section-derived-limits}
for a discussion of derived and homotopy limits and their relationship.
By Derived Categories, Definition \ref{derived-definition-derived-limit}
we have a distinguished
triangle
$$
R\lim K_n \to \prod K_n \to \prod K_n \to R\lim K_n[1]
$$
Taking the long exact sequence of cohomology sheaves we obtain
$$
H^{p - 1}(\prod K_n) \to H^{p - 1}(\prod K_n) \to
H^p(R\lim K_n) \to H^p(\prod K_n) \to H^p(\prod K_n)
$$
Since products are exact by (4) this becomes
$$
\prod H^{p - 1}(K_n) \to \prod H^{p - 1}(K_n) \to
H^p(R\lim K_n) \to \prod H^p(K_n) \to \prod H^p(K_n)
$$
Now we first apply this to the case $K_n = \mathcal{F}_n[0]$
where $(\mathcal{F}_n)$ is as in (6). We conclude that (6) holds.
Next we apply it to $(K_n)$ as in (7) and we conclude (7) holds.
\end{proof}

\end{document}
